% ============================================================
% ВВЕДЕНИЕ
% ============================================================
\section*{ВВЕДЕНИЕ}
\addcontentsline{toc}{section}{ВВЕДЕНИЕ}

\subsection*{Актуальность темы}
\addcontentsline{toc}{subsection}{Актуальность темы}

Дисциплина <<Сопротивление материалов>> входит в базовый цикл инженерного образования и формирует у студентов понимание связи между внешним нагружением конструкции, внутренними силовыми факторами и деформациями~\cite{feodosev, aleksandrov}. Овладение этой дисциплиной является необходимым условием последующего изучения курсов <<Теория упругости>>, <<Строительная механика>>, <<Детали машин>>, <<Проектирование конструкций>> и целого ряда специальных дисциплин инженерного образования~\cite{timoshenko, darkov}. Успешное усвоение основ сопротивления материалов определяет способность будущего инженера к самостоятельному проектированию несущих элементов машин, приборов, зданий и сооружений.

Согласно Федеральному государственному образовательному стандарту высшего образования по направлению 01.03.02 <<Прикладная математика и информатика>>, а также стандартам смежных технических направлений, курс сопротивления материалов относится к базовой части профессионального цикла и предусматривает не менее~144 академических часов аудиторной и самостоятельной работы. Трудоёмкость курса и его положение в учебном плане (обычно 3--4 семестр) делают его одним из наиболее значимых <<фильтров>> инженерного образования: по данным отечественных и зарубежных исследований, от~25\% до~40\% студентов технических специальностей испытывают устойчивые затруднения при первом знакомстве с материалом~\cite{steif2005, philpot2003, makarenko}.

На практике у студентов технических вузов наиболее устойчивые трудности возникают при переходе от аналитических формул и расчётных схем к физическому смыслу результата: почему меняется знак изгибающего момента, где появляется скачок на эпюре поперечных сил, почему максимальный прогиб возникает не в точке приложения нагрузки, и как параметры материала и геометрии сечения влияют на форму деформированной оси балки~\cite{belyaev, pisarenko}. Эти трудности имеют системный характер и связаны не столько с математической сложностью изучаемых уравнений, сколько с необходимостью одновременно удерживать в сознании несколько абстрактных объектов: расчётную схему с приложенными нагрузками, гипотетическое <<рассечение>> балки, распределение внутренних усилий в сечении и результирующую деформированную форму.

Исследования в области инженерной педагогики подтверждают, что активное обучение с использованием интерактивных инструментов значительно повышает уровень усвоения материала по сравнению с традиционными лекционными форматами~\cite{freeman2014, makarenko}. В частности, крупное мета-аналитическое исследование С.~Фримана и соавт.~\cite{freeman2014}, охватившее~225 независимых работ по инженерному и естественнонаучному образованию, показало, что применение методов активного обучения повышает средний балл на экзаменах на~6\% и снижает вероятность получения неудовлетворительной оценки в~1{,}5~раза по сравнению с традиционными лекционными форматами. Аналогичные результаты получены в работах~\cite{makarenko, solovyev, balandina} на выборках российских технических вузов.

В условиях массового перехода высшего образования на цифровые форматы, ускоренного пандемией~COVID-19, потребность в интерактивных обучающих инструментах многократно возросла. Данные Министерства науки и высшего образования Российской Федерации свидетельствуют о том, что доля дистанционных и гибридных форматов обучения в технических вузах в~2020--2024~годах выросла с~12\% до~38\%. В этих условиях качество цифровых учебных материалов напрямую определяет качество инженерной подготовки.

Существующие онлайн-калькуляторы по балкам (sopromatguru.ru, sopromatu.net, sopromat.site, sopromat.xyz) в основном ориентированы на получение конечного результата --- реакций опор, эпюр внутренних силовых факторов, прогибов --- и часто воспринимаются студентами как <<чёрный ящик>>~\cite{philpot2006}: пользователь вводит данные и получает графики без достаточного пояснения причинно-следственных связей. Это снижает учебный эффект, поскольку результат трудно проверить, объяснить и перенести на новые расчётные схемы~\cite{yuan2002}.

В связи с этим актуальной является разработка интерактивной обучающей платформы по основам сопротивления материалов, которая объединяет:
\begin{itemize}
  \item визуализацию эпюр внутренних силовых факторов $M(x)$, $Q(x)$, $N(x)$ и деформированной формы балки;
  \item мгновенный пересчёт при изменении параметров расчётной схемы и нагрузок;
  \item контекстные подсказки, объясняющие физические закономерности;
  \item набор готовых учебных сценариев, приближенных к реальным инженерным ситуациям;
  \item пошаговое решение с демонстрацией метода сечений;
  \item режим тренажёра для самостоятельной проверки знаний;
  \item экспорт результатов в PDF и возможность обмена моделью по ссылке.
\end{itemize}

\subsection*{Обзор опубликованных работ по теме}
\addcontentsline{toc}{subsection}{Обзор опубликованных работ по теме}

Проблематика эффективного преподавания сопротивления материалов и применения интерактивных цифровых инструментов в инженерном образовании изучается с конца~1990-х годов. Обзор опубликованных работ позволяет выделить несколько основных направлений исследований.

\textbf{Направление~1. Диагностика концептуальных трудностей обучения сопромату.} Работы П.\,С.~Стайфа~\cite{steif2005} и Т.\,А.~Филпота~\cite{philpot2006} посвящены систематическому анализу типовых ошибок студентов на начальных этапах изучения курса. Авторами разработан Concept Inventory for Statics --- стандартизированный тест из~27~вопросов, позволяющий количественно измерить глубину понимания фундаментальных понятий. Исследования на выборке более~1000 студентов показали, что традиционные лекционные форматы обеспечивают прирост концептуального понимания лишь на~10--15\%, тогда как методы активного обучения с визуальной поддержкой дают прирост в~30--45\%. Достоинством указанных работ является строгая эмпирическая база; недостатком --- локализация исследований на англоязычной аудитории и отсутствие готовых переносимых инструментов.

\textbf{Направление~2. Разработка интерактивных виртуальных лабораторий.} Работа А.\,К.~Байбулова и соавт.~\cite{baybulov} описывает виртуальную лабораторию по сопромату, реализованную средствами настольного ПО. Авторы демонстрируют эффективность виртуальных стендов при изучении испытаний на растяжение и сжатие. Аналогичный подход применён в работе Е.\,Е.~Истратовой и П.\,В.~Ласточкина~\cite{istratova} для теоретической механики. Достоинство подхода --- погружение в физический эксперимент; недостаток --- привязка к настольному ПО и отсутствие мгновенного перерасчёта параметрических моделей.

\textbf{Направление~3. Компьютерная визуализация как дидактический приём.} В работах И.\,В.~Баландиной~\cite{balandina}, И.\,Ю.~Василенко и И.\,В.~Султановой~\cite{vasilenko} показано, что систематическое применение визуальных образов в техническом образовании повышает как краткосрочное усвоение, так и долгосрочное удержание материала. Авторами обоснован принцип <<наглядности>> как универсальный дидактический принцип цифровой эпохи. Ограничение работ --- отсутствие конкретных программных реализаций, применимых непосредственно к курсу сопромата.

\textbf{Направление~4. Онлайн-платформы для инженерных расчётов.} Ряд зарубежных авторов~\cite{philpot2003, yuan2002, panagiotopoulos2016} разрабатывают веб-инструменты для учебных расчётов в сопротивлении материалов. Работа Т.\,А.~Филпота <<MDSolids>>~\cite{philpot2003} является одним из наиболее известных примеров: это настольное приложение, покрывающее большинство задач базового курса. Недостатком является устаревший интерфейс и отсутствие адаптации к мобильным устройствам. Работа Ю.\,Юаня~\cite{yuan2002} описывает первые попытки переноса учебных расчётов в веб-среду; ограничением служат технические возможности HTML на момент публикации.

\textbf{Направление~5. Современные веб-технологии и педагогические подходы.} Работа Макаренко О.\,В.~\cite{makarenko} и Соловьёва М.\,А. с соавт.~\cite{solovyev} посвящены стратегиям внедрения электронного обучения в российских технических вузах. Авторы обосновывают необходимость создания инструментов с мгновенной обратной связью и адаптивным содержанием. Работа С.\,П.~Пирогова и соавт.~\cite{pirogov} описывает комплекс обучающих программ по теоретической механике, реализованный средствами десктоп-разработки. Общий вывод: существует значительный разрыв между потребностями преподавания и имеющимися инструментами, особенно в части современных веб-технологий.

\textbf{Место настоящей работы в обзоре.} Настоящая выпускная квалификационная работа восполняет указанный пробел, предлагая интерактивную обучающую веб-платформу, построенную на современном клиентском стеке (React~19 + TypeScript~5.9 + Vite), сочетающую принципы активного обучения~\cite{freeman2014}, визуализации~\cite{balandina} и мгновенной обратной связи через пересчёт и перерисовку модели. В отличие от существующих решений, платформа SopromatLab полностью переводит обучение в парадигму непосредственного манипулирования моделью в реальном времени при сохранении педагогически обоснованной структуры материала.

\subsection*{Анализ существующих решений}
\addcontentsline{toc}{subsection}{Анализ существующих решений}

Для оценки текущего состояния рынка обучающих инструментов по сопротивлению материалов был проведён анализ четырёх наиболее популярных русскоязычных онлайн-сервисов.

\textbf{Sopromatguru.ru} позиционируется как облачный сервис для онлайн-расчётов балок, рам и ферм с построением эпюр внутренних усилий. Расчёт балки выполняется с применением математического аппарата метода конечных элементов. Сервис предоставляет широкий охват задач, примеры с разбором, автоматическое построение эпюр. Однако основной акцент сделан на получении результата, а не на обучении через управляемую интерактивность; отсутствуют контекстные подсказки и режим мгновенного пересчёта при изменении параметров.

\textbf{Sopromatu.net} содержит модуль <<Балка-онлайн>>, вычисляющий реакции опор, строящий эпюры, углы поворота и прогибы с выдачей подробно расписанного решения. Имеется банк решённых задач и методические материалы. Тем не менее обучающий эффект ограничен текстовым объяснением: отсутствуют визуальные механики (подсветка участков, объяснение скачков и симметрии), нет единого учебного маршрута и функций совместного использования модели.

\textbf{Sopromat.site} предоставляет калькуляторы для расчёта балок и рам, включая статически неопределимые системы. Платформа воспринимается как инженерный калькулятор, а не учебный тренажёр: отсутствуют встроенные подсказки, учебные сценарии и педагогический слой.

\textbf{Sopromat.xyz} реализует расчёт балки с построением схемы, эпюр, определением прогибов и подбором сечения. Присутствует подробный ход решения и примеры. Основной сценарий использования --- <<построить и получить решение>>, а не обучение через мгновенную визуальную обратную связь.

Сравнительный анализ функциональности существующих решений и разрабатываемой платформы приведён в таблице~\ref{tab:comparison}.

\begin{table}[H]
\caption{Сравнительный анализ существующих решений}\label{tab:comparison}
\small
\begin{tabularx}{\textwidth}{|p{3.5cm}|X|X|X|X|X|}
\hline
\textbf{Критерий} & \textbf{sopromat\-guru.ru} & \textbf{sopro\-matu.net} & \textbf{sopromat\-.site} & \textbf{sopromat\-.xyz} & \textbf{Sopromat\-Lab} \\
\hline
Основной фокус & Расчёт & Расчёт + база задач & Калькуляторы & Расчёт + решение & Обучение \\
\hline
Мгновенный пересчёт & Нет & Нет & Нет & Частично & Да \\
\hline
Контекстные подсказки & Нет & Ограничен. & Нет & Нет & Да (21~условие) \\
\hline
Учебные сценарии & Примеры & Банк задач & Нет & Примеры & Да (5~кейсов) \\
\hline
Тренажёр & Нет & Нет & Нет & Нет & Да \\
\hline
Пошаговое решение & Нет & Текст & Нет & Текст & Интерактив. \\
\hline
Деформации & Частично & Да & Есть & Есть & Ключевой элемент \\
\hline
Share-URL & Нет & Нет & Нет & Нет & Да \\
\hline
Экспорт PDF & Нет & Нет & Нет & Нет & Да \\
\hline
\end{tabularx}
\end{table}

Проведённый анализ показывает, что существующие ресурсы закрывают потребность в быстром расчёте и частично --- в проверке решений. Общий недостаток~--- отсутствие полноценной обучающей интерактивной среды, где студент не только получает эпюры, но и понимает происхождение формы диаграмм и деформированной оси через управляемые изменения параметров и контекстные объяснения~\cite{philpot2003, panagiotopoulos2016}.

\subsection*{Цель и задачи работы}
\addcontentsline{toc}{subsection}{Цель и задачи работы}

\textbf{Цель работы} --- разработка клиентской веб-платформы для интерактивного изучения основ сопротивления материалов с мгновенной визуализацией напряжённо-деформированного состояния конструкций.

Для достижения поставленной цели сформулированы следующие \textbf{задачи}:
\begin{enumerate}
  \item Формализация математических моделей для трёх типов статически определимых конструкций (консольная балка, шарнирно-опёртая балка, балка с вылетом).
  \item Разработка алгоритмов численного расчёта эпюр изгибающих моментов~$M(x)$, поперечных сил~$Q(x)$, продольных сил~$N(x)$ и прогибов~$v(x)$.
  \item Проектирование архитектуры приложения по методологии Feature-Sliced Design.
  \item Реализация интерактивного редактора балок с визуализацией деформированной формы в реальном времени.
  \item Создание образовательных модулей: раздела теории, пошагового решения, контекстных подсказок, тренажёра и учебных сценариев.
  \item Реализация функций экспорта отчётов в PDF и обмена конфигурацией по ссылке.
\end{enumerate}

\subsection*{Научная новизна}
\addcontentsline{toc}{subsection}{Научная новизна}

Научная новизна работы заключается в следующих аспектах:
\begin{enumerate}
  \item \textbf{Архитектурная новизна.} Предложена и реализована архитектура клиентской обучающей платформы на методологии Feature-Sliced Design с применением современного стека React~19 + TypeScript~5.9 + Vite, интегрирующая аналитическое расчётное ядро и образовательный слой (теория, сценарии, тренажёр) в единое интерактивное приложение без серверной части. В отличие от существующих решений, платформа обеспечивает мгновенный отклик при любом изменении параметров модели.

  \item \textbf{Алгоритмическая новизна.} Предложен подход к организации расчётного ядра на основе метода суперпозиции с численным интегрированием методом трапеций при~500 точках дискретизации, обеспечивающий время расчёта менее~5~мс для типовых задач и время полного отклика (пересчёт и перерисовка эпюр) менее~20~мс, что соответствует порогу <<мгновенного>> восприятия~(100~мс)~\cite{flanagan2020}.

  \item \textbf{Педагогическая новизна.} Реализована система контекстных подсказок, содержащая~21~условие, классифицированных по трём группам: физические закономерности, конструктивные особенности и предупреждения. Подсказки активируются автоматически при определённых параметрах модели и объясняют студенту наблюдаемые эффекты (симметрию, экстремумы, перегибы, предупреждения о~чрезмерных деформациях), что позволяет перейти от пассивного наблюдения к активному анализу.
\end{enumerate}

\subsection*{Практическая значимость}
\addcontentsline{toc}{subsection}{Практическая значимость}

Разработанная платформа SopromatLab представляет собой инструмент для самостоятельной подготовки студентов технических специальностей по дисциплине <<Сопротивление материалов>>. Платформа может использоваться:
\begin{itemize}
  \item в качестве демонстрационного средства на лекциях для визуализации эпюр;
  \item для самостоятельной работы студентов с автоматической проверкой результатов;
  \item как тренажёр для подготовки к контрольным работам и экзаменам;
  \item для формирования расчётных отчётов в формате PDF.
\end{itemize}
